% ============================================================
% 01_front.tex — title, abstract, and Introduction (W1)
% Pure fragment: no \documentclass, no \begin{document}.
% \maketitle is issued here; main.tex must not call it again.
%
% Title candidates (T1/T2/T3; chosen = Candidate A):
%   Candidate B: Localize, then Generate: Faithful Interleaved
%                Image-Text Answers from Video
%   Candidate C: Generating Keyframe-Faithful Interleaved
%                Image-Text Answers from Video
% ============================================================

\title{VKIG: Keyframe-Grounded Interleaved Image-Text Generation from Video}
\author{Zizhou Liu\\City University of Hong Kong}
\date{}

\maketitle

\begin{abstract}
Turning moments inside a video into faithful, presentable visual assets is an
increasingly important capability for video platforms, enterprise
documentation, and short-form media production.
However, existing video-language systems either describe video content in text
alone or generate images with no provenance in a source frame, and recent
any-to-any models map video to generated images without localization or
verifiable faithfulness.
To address this, we propose VKIG, a two-agent framework in which an
understanding agent localizes question-relevant evidence and a generation
agent renders images that are newly generated from, and faithful to,
explicitly localized keyframes, coupled through a shared multimodal
retrieval-augmented memory.
Specifically, a training-free two-stage selector combines coarse temporal
windowing by a multimodal language model with object-conditioned in-window
selection, addressing keyframe timing, the bottleneck identified by our
temporal-oracle study; we further formalize the task and construct
VKIG-Bench, a benchmark scoring keyframe consistency and spatio-temporal
localization alongside text and image quality.
Experiments on HC-STVG v2 show that the selector raises joint localization
accuracy to 0.567, surpassing a bridged zero-shot spatio-temporal grounding
baseline at 0.447 on identical clips, and that our faithfulness protocol
exposes generation failures which automatic similarity metrics cannot
reliably detect.
\end{abstract}

\section{Introduction}\label{sec:intro}

Video is now the dominant medium on consumer platforms, and a large share of
practical content work consists of turning moments inside video into
presentable visual assets.
A creator needs a cover image that depicts a real highlight of a clip, an
enterprise converts screen recordings into step-by-step illustrated manuals,
and a short-drama studio produces posters from its episodes.
In each case the desired answer is an image with a few words of text,
traceable to a real moment yet rendered for presentation; faithfulness here
means semantic consistency with the source moment, such as the same subject,
scene, and composition, rather than pixel-level identity recovery.
However, frontier multimodal large language models (MLLMs) can describe video
content yet neither localize the relevant moment precisely nor produce an
accountable image, and text-to-image models generate images with no
provenance in a source video.

Existing routes each cover part of this task.
Retrieval-augmented video question answering, including VideoRAG
\citep{videorag2502} and Video-RAG \citep{videorag2411}, reads long videos
but answers in text only.
Retrieval-augmented image generation methods such as ImageRAG
\citep{imagerag2502} and ORIG \citep{orig2510}, together with interleaved
generators such as M2RAG \citep{m2rag2411} and M2IO-R1 \citep{m2ior12508},
produce image-text output but take text or image input, and the displayed
images are usually retrieved rather than newly generated from a specific
source frame.
Any-to-any models such as X-VILA \citep{xvila2405} and NExT-GPT
\citep{nextgpt2309} can map video input to a generated image, and unified
generators \citep{emu352510,showo22506} natively produce interleaved
image-text; none of them, however, provides explicit spatio-temporal
localization of the evidence or a stated protocol verifying that the
generated image is faithful to a localized source keyframe.
A concurrent preprint \citep{pvtg2607} does generate faithful thumbnails from
video, but it emits a single cover image without an explicit localization
protocol or an interleaved answer.

We introduce VKIG, a two-agent framework for \emph{video-grounded interleaved
image-text generation}.
Given a video and a textual instruction, an \emph{Understand Agent} extracts
the queried target objects, temporally localizes the relevant event, grounds
the objects in the selected keyframe, and drafts the textual answer; a
\emph{Generation Agent} turns each grounded keyframe into a newly generated,
faithfulness-constrained image and lays out the interleaved answer
(Section~\ref{sec:method}).
The two agents are designed to share a multimodal retrieval-augmented memory
with feedback loops between grounding and reselection and between generation
and memory write-back; the reported experiments execute the localization,
grounding, and generation backbone of this design.
The main pipeline is training-free end to end.

A central empirical question is where such a pipeline fails.
A temporal-oracle study shows that keyframe timing, rather than spatial
grounding, is the dominant bottleneck: when the selected keyframe falls
inside the gold temporal window, grounding is accurate, but frame-wise
appearance scoring often misses the window because it fires whenever the
subject is visible.
We therefore design a two-stage selector that adapts the coarse-to-fine
paradigm of long-video question answering \citep{unitime2506,focus2510} to
keyframe selection for faithful generation: an MLLM first localizes a coarse
temporal window over an ordered sparse frame grid, and object-conditioned
scoring then selects the keyframe within that window.
The selector requires no parameter training or adaptation.

To make these properties measurable, we formalize the task and construct
VKIG-Bench (Section~\ref{sec:bench}), to our knowledge the first benchmark
that scores keyframe consistency, that is, whether a generated image
preserves the target object of its source keyframe, and spatio-temporal
localization, in addition to text correctness, image-text alignment, and
image quality.
On the generation side we specify a keyframe-faithfulness protocol with human
and MLLM-judge evaluation, which lets us study how instruction design and
multi-pass editing trade faithfulness against artistry.

On an HC-STVG v2 subset, the two-stage selector raises joint localization
accuracy from 0.433 to 0.567 and outperforms a bridged zero-shot
spatio-temporal video grounding method \citep{zeroshotstvg2509} on identical
clips, and the improvement persists at larger scale
(Section~\ref{sec:exp}).
Our faithfulness study finds that standard automatic similarity metrics
separate wholesale failures well yet lose discriminative power once failures
become expression-level.
Faithfulness proves largely controllable through instruction design and
sequential editing, whereas a unified-model control produces stylized but
unfaithful answers with no localization output.
A real-video prototype runs end to end on a single consumer GPU.

Our contributions are summarized as follows:
\begin{itemize}
\item We propose VKIG, a two-agent framework that couples video understanding
and image generation through a shared multimodal retrieval-augmented memory,
producing interleaved answers in which every image is newly generated from an
explicitly localized keyframe.
\item We design a training-free two-stage keyframe selector that combines
coarse MLLM temporal windowing with object-conditioned in-window selection,
addressing the timing bottleneck identified by our temporal-oracle study and
reaching a joint localization accuracy of 0.567 on HC-STVG v2.
\item We formalize video-grounded interleaved image-text generation and
construct VKIG-Bench, a benchmark scoring keyframe consistency and
spatio-temporal localization alongside text and image quality, together with
a human-anchored keyframe-faithfulness protocol.
\end{itemize}
